\section{Determinación de la constante y covariancia de sus realizaciones}
\label{act21:centralidad-constante}

La constante holográfica de acoplamiento aritmético--geométrico reúne una
determinación estructural y una capacidad de restitución. La primera fija
el registro ordenado que resulta compatible con las incidencias y con la
cronología del estado generado. La segunda establece cómo recuperar ese
registro desde sus representaciones y cómo reconocer las operaciones que
actúan sobre él. La composición de ambas propiedades da contenido preciso
a su función de acoplamiento.

La determinación comienza en las regiones producidas por APP--TRIT--TPK.
Sus bandas de seis trits conservan la posición de cada símbolo, la carga,
el acarreo y la profundidad. El reloj entero compara las escalas ternaria
y decimal; su primer avance nulo, seguido del retorno nonádico, determina
la lectura fase--carga del descriptor. El repertorio terminal y sus
incidencias se conservan con esa procedencia. Las ordenaciones de este
repertorio producen el dominio finito sobre el que actúa la selección
conjunta de la sección~\ref{act21:chap:despliegue-selector-k}.

La condición excepcional restringe ese dominio a los registros cuya cara
de soporte coincide con la cara extraída de las palabras regionales. La
condición temporal compara, dentro de una misma ordenación, las lecturas
comparativa y afín de los dos tramos terminales. De este modo, la posición
de una coordenada y la relación entre dos instantes participan en la misma
decisión. El resultado exacto es
\[
 19\,446\longrightarrow79\longrightarrow1.
\]
El único registro superviviente determina el entero
\[
 N_K=234543140729659824621058914794146601,
 \qquad
 \kappa_{\rm ag}=\frac{N_K}{10^{36}-1}.
\]
La expresión racional conserva el orden de las doce tríadas en la coordenatización
de fase especificada. La unicidad se refiere al dominio terminal
construido, con sus dos condiciones explícitas y sus datos de orientación.
Así queda fijado qué se selecciona, mediante qué relaciones y dentro de
qué conjunto de alternativas.

El censo de las tres terminaciones regionales conserva otra función.
Sus \(196\cdot197\cdot196\) elementos son ternas de palabras de seis
trits compatibles con los cilindros declarados. El perfil de Gram, los
pesos, las distancias y los márgenes orientados determinan la terminación
visible. El censo de \(19\,446\) elementos examina registros ordenados de
doce coordenadas. Ambos reciben información de la misma historia
enriquecida; sus tipos y sus filtros permiten seguir cómo se relacionan
sin confundir una terna visible con el registro completo.

Una vez seleccionado \(K\), la transformación de Hadamard y sus lectores
establecen equivalencias entre representaciones de ese mismo registro.
La selección aporta la determinación entre candidatos; la inversión
aporta la restitución desde una representación. La ley bilateral
proporciona después la identidad de conservación para las isometrías
declaradas, y el archivo conserva la información que permite invertir
sucesivas lecturas. Cada resultado transmite al siguiente un objeto con
dominio, orientación y normalización definidos.

Esta jerarquía precisa también el papel de \(K\) en la ley de
conservación. La identidad bilateral se demuestra para dos isometrías
con dominio y codominio comunes. En las realizaciones del artículo, el
registro seleccionado fija las lecturas aritméticas e incidenciales,
su dirección centrada y las operaciones de reconstrucción que se
componen con aquella identidad. La generalidad del balance y la
determinación del registro son propiedades complementarias: la primera
describe una ley operatoria; la segunda identifica la estructura sobre
la que se calculan las realizaciones aquí expuestas.

La clausura de la constante de estructura fina utiliza las coordenadas
regionales ya generadas y el registro seleccionado, conservando los
acarreos de la operación dodecafásica. Su representación analítica
desarrolla la misma precoordenada con la recurrencia completa. Por la
otra proyección, las incidencias conservadas se transportan a las
realizaciones excepcionales. La procedencia común permite relacionar
valor aritmético, incidencia y memoria sin convertir sus diferentes
lectores en una única magnitud.

Las realizaciones geométricas posteriores amplían esta función. La
dirección centrada de \(K\) determina el flujo diagonal de determinante
uno, preserva los signos de Plücker y el covolumen de las redes
especificadas; el archivo permite recuperarla y controlar sus errores de
lectura. La transformación theta--Mellin conserva la estructura polar en
el dominio demostrado, mientras la diferencia entera registra la
variación espectral. La selección temporal queda así vinculada, mediante
las operaciones ya expuestas, con una estructura que admite evaluación,
transporte y recuperación.

La aportación conjunta del artículo se organiza, por tanto, en una
secuencia demostrativa definida: generación del estado, selección de la
constante, equivalencia de sus representaciones, conservación operatoria
y restitución estable. Esa secuencia determina la posición de cada
teorema y mantiene visibles las condiciones que recibe y las relaciones
que transmite.
